% !TEX root = ../frode-bibelen.tex
% ----------------------------------------------------------------------------
%  Frode-bibelen -- 16 Sjekkliste -- oppgave 2
%  Inkluderes fra frode-bibelen.tex (kompiler alltid hovedfilen).
% ----------------------------------------------------------------------------

% =================================================== 16 SJEKKLISTE 2 ========
\kapittel{Sjekkliste -- oppgave 2}\label{kap:sjekk2}

\subsection*{For hver funksjon du skriver}
\begin{multicols}{2}\raggedright
\boks Lest hele deloppgaven -- og hva som er ferdiglaget\\
\boks Tegnet datastrukturen\\
\boks Ny: \cc{if (gAntall < MAX)} + «Fullt»-melding\\
\boks Andre: \cc{if (gAntall > 0)} før \cc{lesInt(t, 1, gAntall)}\\
\boks Bruker \cc{nr-1} som indeks\\
\boks \cc{malloc} for ny struct \emph{og} for hver \cc{char*}\\
\boks Nullstilt tellere og arrayer i ny struct\\
\boks \cc{gAntall++} etter ny / \cc{gAntall--} etter sletting\\
\boks \cc{free} inni structen \emph{før} \cc{free} av structen\\
\columnbreak
\boks Global funksjon løkker, struct-funksjon gjør jobben\\
\boks \cc{->} for pekere, \cc{.} for struct-variable\\
\boks \cc{strcmp(\dots) == 0} for like tekster\\
\boks \cc{(float)} før deling når snitt skal ha desimaler\\
\boks Fil: \cc{getc} etter \cc{fscanf} før \cc{fgets}\\
\boks Fil: fjernet \cc{'\bs n'} etter \cc{fgets}\\
\boks Fil: \cc{fclose} inni \cc{if (fil)}\\
\boks Skriv-til-fil har \emph{samme} format som les-fra-fil\\
\boks Kommentert det jeg er usikker på
\end{multicols}

\subsection*{Siste 15 minutter}
\boks Alle deloppgaver har \emph{noe} (kode eller pseudokode) \quad \boks Ingen deloppgave står blank\\
\boks Lagret/levert i Inspera
